\documentclass[11pt]{article}

% Packages
\usepackage[margin=1in]{geometry}
\usepackage{amsmath,amssymb,amsthm,mathtools}
\usepackage{hyperref}
\usepackage{booktabs}
\usepackage{graphicx}
\usepackage{natbib}
\usepackage{tikz}
\usetikzlibrary{positioning,arrows.meta,calc}
\usepackage{caption}
\usepackage{multirow}
\usepackage{array}
\usepackage{enumitem}
\usepackage{algorithm}
\usepackage{algpseudocode}

% Hyperref setup
\hypersetup{
    colorlinks=true,
    linkcolor=blue,
    citecolor=blue,
    urlcolor=blue
}

% Theorem environments (matching foundation paper)
\theoremstyle{definition}
\newtheorem{condition}{Condition}
\newtheorem{definition}{Definition}
\newtheorem{example}{Example}
\newtheorem{remark}{Remark}

\theoremstyle{plain}
\newtheorem{theorem}{Theorem}
\newtheorem{lemma}{Lemma}
\newtheorem{proposition}{Proposition}
\newtheorem{corollary}{Corollary}

% Custom commands (matching foundation paper)
\renewcommand{\v}[1]{\boldsymbol{#1}}
\newcommand{\btheta}{\v{\theta}}
\newcommand{\bomega}{\v{\Omega}}
\newcommand{\E}{\mathbb{E}}
\newcommand{\Var}{\mathrm{Var}}
\newcommand{\tr}{\mathrm{tr}}
\newcommand{\rank}{\mathrm{rank}}
\newcommand{\diag}{\mathrm{diag}}
\newcommand{\Loewner}{\preceq}
\newcommand{\LoewnerStrict}{\prec}

% Information-theoretic notation
\newcommand{\FI}{\mathcal{I}}       % Fisher information
\newcommand{\MI}{\mathrm{I}}        % Mutual information
\newcommand{\Ent}{\mathrm{H}}       % Entropy

% Candidate-set matrix
\newcommand{\Ctilde}{\tilde{C}}

\title{Identifiability and Information Loss in Masked Series Systems}

\author{
    Alexander Towell\\
    Southern Illinois University Edwardsville\\
    \href{mailto:lex@metafunctor.com}{lex@metafunctor.com}\\
    \href{https://orcid.org/0000-0001-6443-9897}{ORCID: 0000-0001-6443-9897}
}

\date{February 2026}

\begin{document}

\maketitle

\begin{abstract}
We develop a quantitative theory of identifiability and information loss
for parametric inference in series systems with masked failure causes.
The foundation paper \citep{towell2025masked} establishes binary
(yes/no) identifiability conditions: a parameter vector is identifiable
if and only if the candidate-set structure has sufficient rank.
This paper extends that result in three directions.
First, we introduce the \emph{confounding graph} and show that its
connected components characterize the partially identifiable
sub-parameters, with the number of identifiable parameter groups
determined by a graph coloring argument.
Second, for general parametric families satisfying the C1--C2--C3
conditions, we decompose the Fisher information matrix as
$\FI_{\mathrm{masked}} = \FI_{\mathrm{complete}} - \Delta(\btheta, \pi)$,
where $\Delta$ is a positive semidefinite masking-loss matrix.
For exponential components under uniform masking, both terms admit
closed forms, and we prove that the information loss is monotone
in the masking cardinality via the data processing inequality.
Third, we define an \emph{effective sample size}
$n_{\mathrm{eff}} = n \cdot \lambda_{\min}(\FI_{\mathrm{masked}}) / \lambda_{\min}(\FI_{\mathrm{complete}})$
that translates the information loss into the equivalent number of
complete-data observations, and develop profile-likelihood diagnostics
for detecting near-non-identifiability in practice.
We connect these results to the missing data framework
\citep{little2002} and to optimal diagnostic design via
$D$-optimality.
The theoretical results are validated by a simulation study that
varies the masking probability, number of components, and candidate-set
structure across exponential and Weibull families.
\end{abstract}

\bigskip
\noindent\textbf{Keywords:}
identifiability, Fisher information, masked failure data,
series systems, competing risks, information loss,
effective sample size, diagnostic design

%=============================================================================
\section{Introduction}
\label{sec:introduction}
%=============================================================================

Estimating component-level reliability from series system data is a
central problem in reliability engineering
\citep{barlow1975,agustin2011}.
When the component cause of failure is masked---observed only as a
candidate set rather than an exact identification---the estimation
problem becomes harder.
The foundation paper \citep{towell2025masked} develops a general
likelihood framework under three conditions on the masking mechanism
(C1--C2--C3) and establishes an identifiability theorem: the component
parameter vector $\btheta$ is identifiable if and only if the augmented
candidate-set matrix $\Ctilde$ has full column rank.

That theorem is binary---it answers ``is $\btheta$ identifiable?'' with
yes or no.
But practitioners face a richer set of questions:
\begin{enumerate}[label=(\roman*)]
    \item \textbf{Partial identifiability.}
    When full identifiability fails, which sub-parameters remain
    identifiable?
    \item \textbf{Quantitative degradation.}
    When $\btheta$ is identifiable, how much precision is lost
    relative to complete-data (unmasked) observation?
    \item \textbf{Design guidance.}
    How should the diagnostic procedure be designed to maximize
    information about $\btheta$?
\end{enumerate}
This paper addresses all three questions, developing a theory that
extends the binary identifiability result into a quantitative framework
for information loss and diagnostic design.

The paper is organized in four parts.
Part~I (Sections~\ref{sec:prelim}--\ref{sec:rank-condition}) develops
the structural theory: the confounding graph, its connected components
as the units of partial identifiability, and the matrix-rank condition
for the exponential case.
Part~II (Sections~\ref{sec:fi-decomposition}--\ref{sec:ess}) develops
the information-theoretic story: the Fisher information decomposition,
its closed form for exponential components, the monotone information
loss in masking cardinality, and the effective sample size.
Part~III (Sections~\ref{sec:diagnostics}--\ref{sec:missing-data})
develops the practical tools: near-non-identifiability diagnostics,
optimal diagnostic design, and the connection to missing data theory.
Part~IV (Sections~\ref{sec:simulation}--\ref{sec:discussion}) presents
validation and discussion.

\paragraph{Related work.}
The identifiability of competing risks models has been studied from
several perspectives.
\citet{tsiatis1975} showed that marginal distributions are not
identifiable from competing risks data without independence.
Within the parametric masked-data setting,
\citet{flehinger2002} derived likelihoods for general masked
competing risks and
\citet{lindqvist2023} studied identifiability through phase-type
representations.
Fisher information for masked data has been computed in specific
cases---\citet{towell2026exponential} derives the closed-form FIM
for exponential components under uniform masking---but a general
decomposition into complete-data and masking-loss components has
not been established.
The connection between masking and the missing data framework
of \citet{little2002} has been noted informally
\citep{towell2025masked,craiu2004} but not developed into a
formal fraction-of-missing-information analysis.
Profile-likelihood methods for diagnosing
non-identifiability have been applied in dynamical systems
\citep{raue2009} but not in the masked series system setting.
Optimal experimental design \citep{atkinson2007,pukelsheim2006}
provides the $D$-optimality framework that we apply to the
diagnostic design problem.



%=============================================================================
%  PART I: STRUCTURAL IDENTIFIABILITY
%=============================================================================

%=============================================================================
\section{Preliminaries}
\label{sec:prelim}
%=============================================================================

We adopt the notation and framework of \citet{towell2025masked}
throughout.
A series system consists of $m$ independent components with hazard
functions $h_j(t; \v{\theta}_j)$, reliability functions
$R_j(t; \v{\theta}_j) = \exp\bigl(-\int_0^t h_j(u; \v{\theta}_j)\, du\bigr)$,
and cumulative hazard functions
$H_j(t; \v{\theta}_j) = \int_0^t h_j(u; \v{\theta}_j)\, du$.
The full parameter vector is
$\btheta = (\v{\theta}_1, \ldots, \v{\theta}_m) \in \Theta \subseteq \mathbb{R}^p$.

The system lifetime is $T = \min(T_1, \ldots, T_m)$, and the
component cause of failure is $K = \arg\min_j T_j$.
An observation consists of a system failure time $t_i$, a censoring
indicator $\omega_i$, and a candidate set $c_i \subseteq \{1, \ldots, m\}$.

\begin{condition}[C1--C2--C3; \citealp{towell2025masked}]
\label{cond:c1c2c3}
The masking mechanism satisfies:
\begin{itemize}[nosep]
    \item[\textbf{C1.}] $K_i \in \mathcal{C}_i$ with probability one
        (the candidate set contains the true cause).
    \item[\textbf{C2.}] $\Pr\{\mathcal{C}_i = c_i \mid T_i, K_i = j\}
        = \Pr\{\mathcal{C}_i = c_i \mid T_i, K_i = j'\}$ for all
        $j, j' \in c_i$ (symmetric masking within the candidate set).
    \item[\textbf{C3.}] The masking probabilities, given $(T_i, K_i)$,
        do not depend on~$\btheta$ (masking is ignorable).
\end{itemize}
\end{condition}

Under C1--C2--C3, the likelihood contribution for an exact failure
at time $t_i$ with candidate set $c_i$ is
\begin{equation}
\label{eq:likelihood-contribution}
    L_i(\btheta) \propto \prod_{l=1}^m R_l(t_i; \v{\theta}_l)
    \cdot \sum_{j \in c_i} h_j(t_i; \v{\theta}_j),
\end{equation}
and the log-likelihood for $n$ independent observations (exact
failures and right-censored) is
\begin{equation}
\label{eq:log-likelihood}
    \ell(\btheta) =
    \sum_{i \in \mathcal{E}} \Bigl[
        \sum_{l=1}^m \log R_l(t_i; \v{\theta}_l)
        + \log \sum_{j \in c_i} h_j(t_i; \v{\theta}_j)
    \Bigr]
    + \sum_{i \in \mathcal{R}}
        \sum_{l=1}^m \log R_l(\tau_i; \v{\theta}_l),
\end{equation}
where $\mathcal{E}$ and $\mathcal{R}$ index exact failures and
right-censored observations, respectively.

\begin{definition}[Candidate-set structure]
\label{def:candidate-set-structure}
Let $\mathcal{S} = \{c_1, \ldots, c_s\}$ denote the collection of
candidate sets that occur with positive probability under the masking
mechanism.
The \emph{candidate-set matrix} $C \in \{0,1\}^{|\mathcal{S}| \times m}$
has rows equal to the indicator vectors of the sets in~$\mathcal{S}$.
The \emph{augmented candidate-set matrix} $\Ctilde$ is obtained by
appending the all-ones row $\v{1}^\top$ to~$C$:
\begin{equation}
    \Ctilde = \begin{pmatrix} C \\ \v{1}^\top \end{pmatrix}.
\end{equation}
\end{definition}

\begin{definition}[Candidate-set separability;
\citealp{towell2025masked}]
\label{def:separability}
The masking mechanism is \emph{separating} if, for every pair of
distinct components $j \neq j'$, there exists $c \in \mathcal{S}$
with $j \in c$ and $j' \notin c$.
Components $j$ and $j'$ are \emph{diagnostically confounded} when
$j \in c \iff j' \in c$ for every $c \in \mathcal{S}$.
\end{definition}

We recall the central identifiability result.

\begin{theorem}[Identifiability under C1--C2--C3;
\citealp{towell2025masked}]
\label{thm:identifiability-foundation}
Let each component family be individually identifiable.
Under C1--C2--C3:
\begin{enumerate}[label=(\alph*)]
    \item \textbf{Necessary condition.}
    If $j$ and $j'$ are diagnostically confounded, then $\v{\theta}_j$
    and $\v{\theta}_{j'}$ are not separately identifiable.
    \item \textbf{Sufficient condition.}
    If $\rank(\Ctilde) = m$ and exact-failure events have positive
    probability, then $\btheta$ is identifiable.
\end{enumerate}
\end{theorem}

This paper extends Theorem~\ref{thm:identifiability-foundation}
from a binary verdict to a quantitative characterization.


%=============================================================================
\section{Graph-Theoretic Characterization of Partial Identifiability}
\label{sec:graph}
%=============================================================================

When identifiability fails, some sub-parameters may still be
recoverable.
We formalize this through the \emph{confounding graph}.

\begin{definition}[Confounding graph]
\label{def:confounding-graph}
The \emph{confounding graph} $G = (V, E)$ has vertex set
$V = \{1, \ldots, m\}$ (the components) and edge set
\begin{equation}
    E = \bigl\{(j, j') : j \text{ and } j' \text{ are diagnostically confounded}\bigr\}.
\end{equation}
That is, $(j, j') \in E$ if and only if $j \in c \iff j' \in c$
for every $c \in \mathcal{S}$.
\end{definition}

Diagnostic confounding is an equivalence relation on components: if
$j$ is confounded with $j'$ and $j'$ is confounded with $j''$, then
$j$ is confounded with $j''$ (they share identical column patterns
in~$C$).
The connected components of $G$ are therefore the equivalence classes
of this relation.

\begin{definition}[Super-component]
\label{def:super-component}
A \emph{super-component} is a connected component of $G$.
Let $\mathcal{P} = \{P_1, \ldots, P_q\}$ denote the partition of
$\{1, \ldots, m\}$ into super-components.
For each super-component $P_k$, the \emph{aggregate hazard} is
\begin{equation}
    h_{P_k}(t; \btheta) = \sum_{j \in P_k} h_j(t; \v{\theta}_j).
\end{equation}
\end{definition}

\begin{theorem}[Partial identifiability via the confounding graph]
\label{thm:partial-identifiability}
Under C1--C2--C3, with each component family individually
identifiable:
\begin{enumerate}[label=(\alph*)]
    \item The aggregate hazard $h_{P_k}$ is identifiable for each
    super-component $P_k$ if and only if the augmented candidate-set
    matrix, reduced to have one column per super-component
    (where the column for $P_k$ is the common column shared by all
    $j \in P_k$), has full column rank~$q$.
    \item The individual parameter vectors $\v{\theta}_j$ for $j \in P_k$
    with $|P_k| \geq 2$ are not separately identifiable from the
    observed-data likelihood.
    The likelihood depends on these parameters only through
    $h_{P_k}$.
    \item The number of identifiable super-component parameters is at
    most $q \cdot d$, where $d$ is the dimension of each component's
    parameter vector (assuming a common parametric family).
\end{enumerate}
\end{theorem}

\begin{proof}
\emph{Part~(a).}
Define the reduced candidate-set matrix $C^* \in \{0,1\}^{|\mathcal{S}| \times q}$
by collapsing the columns within each super-component (which are
identical by definition of diagnostic confounding).
Let $\Ctilde^*$ be $C^*$ augmented with the all-ones row.
Each observation's likelihood contribution depends on components in
$P_k$ only through $h_{P_k}$: by
Theorem~\ref{thm:identifiability-foundation}(a), components within
$P_k$ always appear together in every candidate set, so their hazards
always sum.
The problem reduces to identifiability of $q$ ``components'' with
hazards $h_{P_1}, \ldots, h_{P_q}$ and candidate-set matrix $C^*$.
By Theorem~\ref{thm:identifiability-foundation}(b), identifiability
holds when $\rank(\Ctilde^*) = q$.

\emph{Part~(b).}
This is a direct consequence of
Theorem~\ref{thm:identifiability-foundation}(a): within $P_k$,
every pair is diagnostically confounded, so any reparametrization
preserving $h_{P_k}$ leaves the likelihood unchanged.

\emph{Part~(c).}
Each identifiable super-component contributes at most $d$ parameters
to the identifiable part of $\btheta$, and there are $q$ such
super-components.
For super-components with $|P_k| = 1$, all $d$ parameters of
the single member are identifiable.
For super-components with $|P_k| \geq 2$, only the parameters of
the aggregate hazard $h_{P_k}$ are identifiable, which is at most
$d$ parameters (corresponding to fitting a single member of the
parametric family to the aggregate).
\end{proof}

\begin{example}[Hierarchical diagnostics]
\label{ex:hierarchical}
Consider a 5-component system where diagnostics can distinguish
components within subsystems $\{1, 2\}$, $\{3, 4\}$, and $\{5\}$,
but cannot distinguish between subsystems when multiple subsystems
are implicated.
Suppose $\mathcal{S} = \bigl\{\{1,2,3,4\}, \{1,2,5\},
\{3,4,5\}, \{1\}, \{2\}, \{3\}, \{4\}, \{5\}\bigr\}$.
Every pair of components is separated by some candidate set
(e.g., $\{1\}$ separates 1 from all others), so $G$ has no edges,
$q = m = 5$, and all parameters are identifiable.

Now remove the singleton candidate sets:
$\mathcal{S}' = \bigl\{\{1,2,3,4\}, \{1,2,5\}, \{3,4,5\}\bigr\}$.
Components 1 and 2 always co-occur, as do 3 and 4.
The confounding graph has edges $(1,2)$ and $(3,4)$, giving
super-components $P_1 = \{1,2\}$, $P_2 = \{3,4\}$, $P_3 = \{5\}$.
Only $q = 3$ aggregate hazards are identifiable.
\end{example}

\begin{remark}[Remediation strategies]
\label{rem:remediation}
When the confounding graph has non-trivial components, the
practitioner has three options, as described in
\citet{towell2025masked}:
(i)~collapse each super-component into a single composite component
with hazard $h_{P_k}$;
(ii)~impose equality constraints
$\v{\theta}_j = \v{\theta}_{j'}$ for $j, j' \in P_k$ (appropriate
when the components are physically identical); or
(iii)~introduce informative priors to regularize the non-identifiable
directions.
The confounding graph precisely identifies which components require
such treatment.
\end{remark}


%=============================================================================
\section{Candidate-Set Matrix Rank: The Exponential Case}
\label{sec:rank-condition}
%=============================================================================

For exponential components with hazard $h_j(t; \lambda_j) = \lambda_j$
(constant), the identifiability condition simplifies to a linear
algebra problem.
This section provides a self-contained treatment.

\subsection{The exponential likelihood}

Under C1--C2--C3 with exponential components, the log-likelihood
for exact failures and right-censored observations is
\begin{equation}
\label{eq:expo-loglik}
    \ell(\v{\lambda})
    = \sum_{i \in \mathcal{E}} \Bigl[
        -\Lambda \, t_i + \log \sum_{j \in c_i} \lambda_j
    \Bigr]
    - \Lambda \sum_{i \in \mathcal{R}} \tau_i,
\end{equation}
where $\Lambda = \sum_{j=1}^m \lambda_j$ is the total system rate
and $\v{\lambda} = (\lambda_1, \ldots, \lambda_m)^\top \in (0,\infty)^m$.

The key structural feature is that the exponential hazard is constant,
so the second term $\log \sum_{j \in c_i} \lambda_j$ depends on the
candidate set $c_i$ but not on time.
This means the information about $\v{\lambda}$ from the candidate-set
structure is purely combinatorial.

\subsection{The rank condition}

\begin{theorem}[Exponential identifiability via matrix rank]
\label{thm:expo-rank}
For exponential components under C1--C2--C3, the parameter vector
$\v{\lambda}$ is identifiable if and only if
$\rank(\Ctilde) = m$.
\end{theorem}

\begin{proof}
\emph{Sufficiency} is
Theorem~\ref{thm:identifiability-foundation}(b), since exponential
hazards are individually identifiable ($\lambda_j = \lambda_j'$
whenever $h_j \equiv h_j'$).

\emph{Necessity.}
Suppose $\rank(\Ctilde) < m$.
Then there exists a nonzero vector $\v{g} \in \mathbb{R}^m$ with
$\Ctilde \v{g} = \v{0}$.
In particular, $\v{1}^\top \v{g} = 0$ (from the last row) and
$\sum_{j \in c} g_j = 0$ for each $c \in \mathcal{S}$ (from the
rows of $C$).
For any $\epsilon > 0$ sufficiently small, define
$\v{\lambda}' = \v{\lambda} + \epsilon \v{g}$.
Then $\v{\lambda}' \in (0,\infty)^m$ for small enough $\epsilon$,
$\Lambda' = \Lambda$ (since $\v{1}^\top \v{g} = 0$), and
\begin{equation}
    \sum_{j \in c} \lambda_j' = \sum_{j \in c} \lambda_j
        + \epsilon \sum_{j \in c} g_j
    = \sum_{j \in c} \lambda_j
\end{equation}
for every $c \in \mathcal{S}$.
Thus $\ell(\v{\lambda}') = \ell(\v{\lambda})$ for all data sets,
so $\v{\lambda}$ and $\v{\lambda}'$ are not distinguishable.
\end{proof}

\begin{remark}[Exponential vs.\ general families]
\label{rem:expo-vs-general}
For exponential components, the rank condition is both necessary and
sufficient because the hazards $h_j = \lambda_j$ are constant, so
$g_j(t) = \lambda_j - \lambda_j'$ is constant and the equation
$\Ctilde \v{g} = \v{0}$ is independent of $t$.
For families with time-varying hazards (e.g., Weibull with distinct
shapes), the equation $\sum_{j \in c} g_j(t) = 0$ for all $t > 0$
can force $g_j \equiv 0$ even when $\rank(\Ctilde) < m$, because
the functional diversity of the hazards provides additional
constraints.
See \citet[Remark~4]{towell2025masked} for the precise statement.
\end{remark}

\begin{corollary}[Uniform masking identifiability]
\label{cor:uniform-masking}
Under uniform masking with cardinality $w < m$, every candidate set
of size $w$ occurs with positive probability.
Since the collection of all $\binom{m}{w}$ subsets of size $w$
generates a matrix $C$ whose augmentation $\Ctilde$ has full column
rank $m$ (by a counting argument: any vector in the null space must
have all entries equal, and $w < m$ ensures this is ruled out by the
all-ones row), $\v{\lambda}$ is identifiable.
At $w = m$, there is a single candidate set $\{1, \ldots, m\}$,
$\rank(\Ctilde) = 1$, and only $\Lambda$ is identifiable.
\end{corollary}



%=============================================================================
%  PART II: INFORMATION LOSS
%=============================================================================

%=============================================================================
\section{Fisher Information Decomposition}
\label{sec:fi-decomposition}
%=============================================================================

We now move from the qualitative question (is $\btheta$ identifiable?)
to the quantitative question (how precisely can $\btheta$ be estimated?).
The key tool is the Fisher information matrix.

\subsection{Complete-data and masked-data information}

Consider two observation regimes for the same series system:
\begin{enumerate}[nosep]
    \item \textbf{Complete data:} Observe $(T_i, K_i)$---the system
    failure time and the exact component cause.
    \item \textbf{Masked data:} Observe $(T_i, \mathcal{C}_i)$---the
    system failure time and the candidate set (under C1--C2--C3).
\end{enumerate}

Let $\FI_{\mathrm{complete}}(\btheta)$ denote the per-observation
Fisher information matrix under complete data, and
$\FI_{\mathrm{masked}}(\btheta)$ denote the per-observation Fisher
information under masked data.

\begin{theorem}[Fisher information decomposition]
\label{thm:fi-decomposition}
Under C1--C2--C3, assuming standard regularity conditions for the
interchange of differentiation and integration:
\begin{equation}
\label{eq:fi-decomposition}
    \FI_{\mathrm{masked}}(\btheta)
    = \FI_{\mathrm{complete}}(\btheta) - \Delta(\btheta, \pi),
\end{equation}
where $\Delta(\btheta, \pi) \succeq 0$ (positive semidefinite) is
the \emph{masking-loss matrix} that depends on both the component
parameters $\btheta$ and the masking mechanism $\pi$.
\end{theorem}

\begin{proof}
The complete data $(T_i, K_i)$ is a sufficient reduction of the
``full'' data $(T_i, K_i, \mathcal{C}_i)$ for $\btheta$ (since the
candidate set, under C3, carries no information about $\btheta$
beyond what $K_i$ provides).
The masked data $(T_i, \mathcal{C}_i)$ is a further reduction from
$(T_i, K_i)$---the exact cause $K_i$ is replaced by the coarser
candidate set $\mathcal{C}_i$.

By the information inequality for sufficient statistics
\citep[Ch.~7]{schervish1995}, any reduction of the data can only
lose Fisher information:
\begin{equation}
    \FI_{\mathrm{masked}}(\btheta) \Loewner
    \FI_{\mathrm{complete}}(\btheta)
\end{equation}
in the Loewner (positive semidefinite) ordering.
The masking-loss matrix $\Delta = \FI_{\mathrm{complete}} - \FI_{\mathrm{masked}}$
is therefore positive semidefinite.

More constructively, we can derive this via the law of total Fisher
information.
Let $\v{S}(\btheta) = \nabla_\btheta \log f(T_i, \mathcal{C}_i; \btheta)$
be the masked-data score and
$\v{S}^*(\btheta) = \nabla_\btheta \log f(T_i, K_i; \btheta)$
be the complete-data score.
Then
\begin{equation}
    \FI_{\mathrm{complete}}(\btheta)
    = \E[\v{S}^* {\v{S}^*}^\top]
    = \E\bigl[\E[\v{S}^* {\v{S}^*}^\top \mid T_i, \mathcal{C}_i]\bigr]
\end{equation}
and
\begin{equation}
    \FI_{\mathrm{masked}}(\btheta)
    = \E[\v{S} \v{S}^\top]
    = \E\bigl[\E[\v{S}^* \mid T_i, \mathcal{C}_i]
        \E[\v{S}^* \mid T_i, \mathcal{C}_i]^\top\bigr].
\end{equation}
By the conditional variance decomposition,
\begin{align}
    \Delta(\btheta, \pi)
    &= \FI_{\mathrm{complete}} - \FI_{\mathrm{masked}} \notag \\
    &= \E\bigl[\Var(\v{S}^* \mid T_i, \mathcal{C}_i)\bigr]
    \succeq 0,
\end{align}
since the expectation of a (conditional) covariance matrix is positive
semidefinite.
\end{proof}

\begin{remark}[Interpretation]
\label{rem:decomposition-interpretation}
The decomposition \eqref{eq:fi-decomposition} has a natural
interpretation: $\FI_{\mathrm{complete}}$ is the total information
available from the series system experiment,
$\FI_{\mathrm{masked}}$ is the information retained after masking,
and $\Delta$ is the information destroyed by the diagnostic's inability
to identify the exact failure cause.
The masking-loss matrix $\Delta$ captures how much the score function
fluctuates \emph{within} a given candidate set---if the candidate set
leaves substantial ambiguity about which component failed, $\Delta$ is
large.
When $\mathcal{C}_i = \{K_i\}$ (singleton candidate sets, no masking),
$\Delta = 0$ and all information is retained.
When $\mathcal{C}_i = \{1, \ldots, m\}$ (complete masking), $\Delta$
is maximal.
\end{remark}

\subsection{Loewner ordering of information}

\begin{corollary}[Information ordering]
\label{cor:loewner-ordering}
For any scalar function $g(\btheta)$, the asymptotic variance of the
MLE $\hat{g}$ satisfies
\begin{equation}
    \Var_{\mathrm{masked}}(\hat{g})
    \geq \Var_{\mathrm{complete}}(\hat{g}),
\end{equation}
with equality if and only if
$\nabla g^\top \Delta \, \nabla g = 0$---that is, the direction
$\nabla g$ lies in the null space of the masking-loss matrix.
\end{corollary}

\begin{proof}
By the Cram\'er--Rao bound, the asymptotic variance of $\hat{g}$
under regime $r \in \{\text{masked}, \text{complete}\}$ is at least
$(\nabla g)^\top \FI_r^{-1} \nabla g$.
Since $\FI_{\mathrm{masked}} \Loewner \FI_{\mathrm{complete}}$,
we have $\FI_{\mathrm{masked}}^{-1} \succeq \FI_{\mathrm{complete}}^{-1}$
(the Loewner order reverses under inversion for positive definite
matrices), yielding the result.
\end{proof}


%=============================================================================
\section{Exponential Case: Closed-Form Information Loss}
\label{sec:expo-info}
%=============================================================================

We now derive explicit expressions for both $\FI_{\mathrm{complete}}$
and $\FI_{\mathrm{masked}}$ when the components are exponentially
distributed.
This section is self-contained; it provides the essential results
from \citet{towell2026exponential} re-derived within the present
framework.

\subsection{Complete-data Fisher information}

Under complete data, we observe $(T_i, K_i)$ where $T_i$ is the
system failure time and $K_i \in \{1, \ldots, m\}$ is the exact
failure cause.
The joint density is
$f(t, k; \v{\lambda}) = \lambda_k \exp(-\Lambda t)$,
where $\Lambda = \sum_{j=1}^m \lambda_j$.

\begin{proposition}[Complete-data FIM, exponential]
\label{prop:complete-fim}
For exponential components with parameter $\v{\lambda} \in (0,\infty)^m$,
the per-observation complete-data Fisher information matrix is
\begin{equation}
\label{eq:complete-fim}
    [\FI_{\mathrm{complete}}(\v{\lambda})]_{jk}
    = \frac{\delta_{jk}}{\lambda_j \Lambda}
    \quad \text{for } j, k = 1, \ldots, m,
\end{equation}
where $\delta_{jk}$ is the Kronecker delta.
In matrix form:
\begin{equation}
    \FI_{\mathrm{complete}}(\v{\lambda})
    = \frac{1}{\Lambda} \diag\!\left(
        \frac{1}{\lambda_1}, \ldots, \frac{1}{\lambda_m}
    \right).
\end{equation}
\end{proposition}

\begin{proof}
The log-density is
$\log f(t, k; \v{\lambda}) = \log \lambda_k - \Lambda \, t$.
Differentiating:
\begin{equation}
    \frac{\partial \log f}{\partial \lambda_j}
    = \frac{\delta_{jk}}{\lambda_j} - t.
\end{equation}
Taking the second derivative:
\begin{equation}
    \frac{\partial^2 \log f}{\partial \lambda_j \partial \lambda_l}
    = -\frac{\delta_{jl} \delta_{jk}}{\lambda_j^2}.
\end{equation}
Now compute the Fisher information:
\begin{align}
    [\FI_{\mathrm{complete}}]_{jl}
    &= -\E\!\left[\frac{\partial^2 \log f}{\partial \lambda_j
        \partial \lambda_l}\right]
    = \sum_{k=1}^m \int_0^\infty \frac{\delta_{jl} \delta_{jk}}{\lambda_j^2}
        \cdot \lambda_k e^{-\Lambda t}\, dt \\
    &= \frac{\delta_{jl}}{\lambda_j^2} \cdot \lambda_j
        \cdot \frac{1}{\Lambda}
    = \frac{\delta_{jl}}{\lambda_j \Lambda}.
\end{align}
The integral $\int_0^\infty e^{-\Lambda t}\, dt = 1/\Lambda$, and
the sum over $k$ picks out $k = j$ via $\delta_{jk}$.
\end{proof}

\subsection{Masked-data Fisher information under uniform masking}

Under uniform masking with cardinality $w$, the per-observation
masked-data Fisher information was derived by
\citet{towell2026exponential}.
We re-derive it here for self-containment.

\begin{proposition}[Masked-data FIM, exponential;
cf.~\citealp{towell2026exponential}]
\label{prop:masked-fim}
Under uniform masking with cardinality $w$, the per-observation
Fisher information for exponential components is
\begin{equation}
\label{eq:masked-fim}
    [\FI_{\mathrm{masked}}(\v{\lambda} \mid w)]_{jk}
    = \frac{1}{\binom{m-1}{w-1} \Lambda}
    \sum_{\substack{c \subseteq \{1,\ldots,m\} \\ |c| = w}}
    \frac{\mathbf{1}_{j \in c} \, \mathbf{1}_{k \in c}}
         {\sum_{l \in c} \lambda_l}.
\end{equation}
\end{proposition}

\begin{proof}
Under uniform masking, the joint density of $(T_i, \mathcal{C}_i)$ is
\begin{equation}
    f(t, c; \v{\lambda}) = \frac{1}{\binom{m-1}{w-1}}
    \left(\sum_{j \in c} \lambda_j\right) e^{-\Lambda t},
    \quad |c| = w.
\end{equation}
The log-density is
$\log f = \log(\sum_{j \in c} \lambda_j) - \Lambda t + \text{const}$.
The second partial derivative is
\begin{equation}
    \frac{\partial^2 \log f}{\partial \lambda_j \partial \lambda_k}
    = -\frac{\mathbf{1}_{j \in c} \, \mathbf{1}_{k \in c}}
           {(\sum_{l \in c} \lambda_l)^2}.
\end{equation}
Taking the expectation over $(T, \mathcal{C})$:
\begin{align}
    [\FI_{\mathrm{masked}}]_{jk}
    &= \sum_{c: |c|=w} \int_0^\infty
        \frac{\mathbf{1}_{j \in c} \, \mathbf{1}_{k \in c}}
             {(\sum_l \lambda_l)_{l \in c}^2}
        \cdot \frac{(\sum_l \lambda_l)_{l \in c}}{\binom{m-1}{w-1}}
        e^{-\Lambda t}\, dt \\
    &= \frac{1}{\binom{m-1}{w-1} \Lambda}
        \sum_{c: |c|=w}
        \frac{\mathbf{1}_{j \in c} \, \mathbf{1}_{k \in c}}
             {\sum_{l \in c} \lambda_l}.
\end{align}
\end{proof}

\subsection{Closed-form masking-loss matrix}

\begin{theorem}[Masking-loss matrix, exponential]
\label{thm:expo-delta}
For exponential components under uniform masking with cardinality $w$,
the masking-loss matrix is
\begin{equation}
\label{eq:expo-delta}
    \Delta(\v{\lambda}, w)
    = \FI_{\mathrm{complete}}(\v{\lambda})
    - \FI_{\mathrm{masked}}(\v{\lambda} \mid w),
\end{equation}
with entries
\begin{equation}
    [\Delta(\v{\lambda}, w)]_{jk}
    = \frac{\delta_{jk}}{\lambda_j \Lambda}
    - \frac{1}{\binom{m-1}{w-1} \Lambda}
    \sum_{c: |c|=w}
    \frac{\mathbf{1}_{j \in c} \, \mathbf{1}_{k \in c}}
         {\sum_{l \in c} \lambda_l}.
\end{equation}
This matrix is positive semidefinite with $\Delta = 0$ when $w = 1$
(no masking) and $\rank(\Delta) = m - 1$ when $w = m$ (complete masking,
where only $\Lambda$ is estimable).
\end{theorem}

\begin{proof}
The formula follows by subtracting
\eqref{eq:masked-fim} from \eqref{eq:complete-fim}.
Positive semidefiniteness is guaranteed by
Theorem~\ref{thm:fi-decomposition}.

When $w = 1$, the candidate sets are singletons $\{j\}$.
The sum in \eqref{eq:masked-fim} over $c = \{j\}$ gives
$[\FI_{\mathrm{masked}}]_{jk} = \delta_{jk} / (\lambda_j \Lambda)$,
matching $\FI_{\mathrm{complete}}$, so $\Delta = 0$.

When $w = m$, there is a single candidate set $\{1, \ldots, m\}$, so
$[\FI_{\mathrm{masked}}]_{jk} = 1/(\Lambda^2)$ for all $j, k$.
This is a rank-1 matrix: $\FI_{\mathrm{masked}} = \v{1}\v{1}^\top / \Lambda^2$.
The masking-loss matrix
$\Delta = \frac{1}{\Lambda}\diag(\lambda_j^{-1}) - \frac{1}{\Lambda^2}\v{1}\v{1}^\top$
has rank $m - 1$ (its null space is spanned by $\v{\lambda}$, since
$\Delta \v{\lambda} = \v{1}/\Lambda - \v{1}/\Lambda = \v{0}$).
\end{proof}

\begin{example}[Three-component symmetric system]
\label{ex:three-component}
For $m = 3$, $\lambda_1 = \lambda_2 = \lambda_3 = \lambda$,
and $w = 2$:
\begin{equation}
    \FI_{\mathrm{complete}}
    = \frac{1}{3\lambda^2}
    \begin{pmatrix} 1 & 0 & 0 \\ 0 & 1 & 0 \\ 0 & 0 & 1 \end{pmatrix},
    \quad
    \FI_{\mathrm{masked}}
    = \frac{1}{6\lambda^2}
    \begin{pmatrix} 2 & 1 & 1 \\ 1 & 2 & 1 \\ 1 & 1 & 2 \end{pmatrix}.
\end{equation}
The eigenvalues of $\FI_{\mathrm{complete}}$ are all $1/(3\lambda^2)$.
The eigenvalues of $\FI_{\mathrm{masked}}$ are
$4/(6\lambda^2) = 2/(3\lambda^2)$ (multiplicity 1, eigenvector
$\v{1}$) and
$1/(6\lambda^2)$ (multiplicity 2).
The information about the total rate $\Lambda$ is actually
\emph{higher} under masking (the $\v{1}$ direction), because masking
pools information about $\Lambda$ from all candidate sets.
The information about individual rate \emph{differences} (the two
degenerate directions) is reduced by a factor of
$\frac{1/(6\lambda^2)}{1/(3\lambda^2)} = 1/2$.

The masking-loss matrix has eigenvalues 0 (for the $\v{1}$ direction)
and $1/(6\lambda^2)$ (multiplicity 2), confirming that masking
destroys information about \emph{which component is which} but
preserves information about the total system hazard.
\end{example}


%=============================================================================
\section{Monotone Information Loss}
\label{sec:monotone}
%=============================================================================

Intuitively, more masking should mean more information loss.
We formalize this as a monotonicity result in the masking
cardinality~$w$.

\subsection{Mutual information and the data processing inequality}

Before addressing Fisher information monotonicity, we establish the
corresponding result for mutual information (a purely combinatorial
quantity independent of $\btheta$), then connect it to Fisher information.

\begin{definition}[Diagnostic informativeness]
\label{def:diagnostic-info}
Under uniform masking with cardinality $w$, the \emph{diagnostic mutual
information} is
$\MI(K; \mathcal{C}_w) = \Ent(K) - \Ent(K \mid \mathcal{C}_w)$,
measuring how much the candidate set reveals about the failure cause.
\end{definition}

\begin{proposition}[Monotone mutual information loss;
cf.~\citealp{towell2026exponential}]
\label{prop:mi-monotone}
Under uniform masking:
\begin{enumerate}[label=(\alph*)]
    \item The Markov chain $K \to \mathcal{C}_w \to \mathcal{C}_{w+1}$
    holds for $1 \leq w \leq m-1$.
    \item $\MI(K; \mathcal{C}_w)$ is strictly decreasing in $w$ for
    $1 \leq w \leq m$.
    \item Under equal rates ($\lambda_1 = \cdots = \lambda_m$),
    $\MI(K; \mathcal{C}_w) = \ln(m/w)$.
\end{enumerate}
\end{proposition}

\begin{proof}
\emph{Part~(a).}
Given $\mathcal{C}_w$ (a uniform random size-$w$ set containing $K$),
construct $\mathcal{C}_{w+1} = \mathcal{C}_w \cup \{J\}$ where $J$
is drawn uniformly from $\{1, \ldots, m\} \setminus \mathcal{C}_w$.
Conditional on $\mathcal{C}_w$, $\mathcal{C}_{w+1}$ is determined by
the random addition of one element from outside $\mathcal{C}_w$, which
does not depend on $K$ given $\mathcal{C}_w$ (since $K \in \mathcal{C}_w$
by C1).
Hence $K \perp \mathcal{C}_{w+1} \mid \mathcal{C}_w$, establishing
the Markov chain.

We verify that $\mathcal{C}_{w+1}$ has the correct marginal
distribution.
For a fixed size-$(w+1)$ set $c'$ containing $K = k$, the number
of size-$w$ subsets of $c'$ containing $k$ is $\binom{w}{w-1} = w$.
Each such subset arises with probability $1/\binom{m-1}{w-1}$, and
the random element $J$ is chosen from $m - w$ options, so
\begin{equation}
    \Pr\{\mathcal{C}_{w+1} = c' \mid K = k\}
    = \frac{w}{\binom{m-1}{w-1}(m-w)}
    = \frac{1}{\binom{m-1}{w}},
\end{equation}
which is the uniform masking distribution with cardinality $w+1$.

\emph{Part~(b).}
By the data processing inequality
\citep{coverandthomas2006}, the Markov chain
$K \to \mathcal{C}_w \to \mathcal{C}_{w+1}$ implies
$\MI(K; \mathcal{C}_{w+1}) \leq \MI(K; \mathcal{C}_w)$.
Strict inequality holds because the map
$\mathcal{C}_w \mapsto \mathcal{C}_{w+1}$ is not invertible (it
adds an element, losing information about which element was original).

\emph{Part~(c).}
Under equal rates, $K$ is uniform on $\{1, \ldots, m\}$, so
$\Ent(K) = \ln m$.
Given $K = k$ and $\mathcal{C}_w = c$ with $k \in c$,
$\Pr\{K = k \mid \mathcal{C}_w = c\} = 1/w$ (by symmetry of both
the prior and the masking), so $\Ent(K \mid \mathcal{C}_w) = \ln w$.
Therefore $\MI(K; \mathcal{C}_w) = \ln m - \ln w = \ln(m/w)$.
\end{proof}

\subsection{Fisher information monotonicity}

\begin{theorem}[Monotone Fisher information loss]
\label{thm:fi-monotone}
Under uniform masking with exponential components:
\begin{equation}
    \FI_{\mathrm{masked}}(\v{\lambda} \mid w+1)
    \Loewner
    \FI_{\mathrm{masked}}(\v{\lambda} \mid w)
    \Loewner
    \FI_{\mathrm{complete}}(\v{\lambda})
\end{equation}
for $1 \leq w \leq m - 1$.
Equivalently, $\Delta(\v{\lambda}, w)$ is increasing in $w$ in the
Loewner order.
\end{theorem}

\begin{proof}
The Markov chain $K \to \mathcal{C}_w \to \mathcal{C}_{w+1}$
established in Proposition~\ref{prop:mi-monotone}(a) gives
$(T, \mathcal{C}_{w+1})$ as a coarsening of $(T, \mathcal{C}_w)$.
By the Fisher information data processing inequality---which states
that any statistic computed from the data can only lose Fisher
information about the parameter \citep{zamir1998}---we have
$\FI_{\mathrm{masked}}(\v{\lambda} \mid w+1)
\Loewner
\FI_{\mathrm{masked}}(\v{\lambda} \mid w)$.

The right inequality
$\FI_{\mathrm{masked}}(\v{\lambda} \mid w)
\Loewner \FI_{\mathrm{complete}}(\v{\lambda})$
is the special case $w' = 1$ (where
$\FI_{\mathrm{masked}}(\v{\lambda} \mid 1) = \FI_{\mathrm{complete}}(\v{\lambda})$).
\end{proof}

\begin{remark}[Maximum-entropy characterization]
\label{rem:max-entropy}
Among all masking models satisfying C1 and C2 with fixed candidate-set
size $w$, \emph{uniform masking maximizes} the conditional entropy
$\Ent(\mathcal{C} \mid K)$
\citep{towell2026exponential}.
Combined with the monotonicity result, this identifies uniform masking
at $w = m - 1$ as the maximally uninformative identifiable scenario
under C2: it maximizes entropy and minimizes mutual information among
all identifiable uniform masking configurations.
This makes uniform masking with $w = m - 1$ a natural \emph{worst-case
benchmark} for practitioners assessing the impact of masking.
\end{remark}


%=============================================================================
\section{Effective Sample Size}
\label{sec:ess}
%=============================================================================

The Fisher information decomposition provides a matrix-valued measure
of information loss.
For practical communication, we distill this into a scalar:
the \emph{effective sample size}.

\begin{definition}[Effective sample size]
\label{def:ess}
Given $n$ masked observations, the \emph{effective sample size}
for estimating $\btheta$ is
\begin{equation}
\label{eq:ess-def}
    n_{\mathrm{eff}}
    = n \cdot r(\btheta, \pi),
\end{equation}
where $r(\btheta, \pi)$ is the \emph{information retention ratio}:
\begin{equation}
\label{eq:irr}
    r(\btheta, \pi)
    = \frac{\lambda_{\min}\bigl(\FI_{\mathrm{masked}}(\btheta)\bigr)}
           {\lambda_{\min}\bigl(\FI_{\mathrm{complete}}(\btheta)\bigr)}.
\end{equation}
Here $\lambda_{\min}(\cdot)$ denotes the smallest eigenvalue.
\end{definition}

The effective sample size answers the question: ``how many
complete-data observations would provide the same worst-case
precision as my $n$ masked observations?''

\begin{remark}[Alternative definitions]
\label{rem:ess-alternatives}
The minimum-eigenvalue ratio in \eqref{eq:irr} captures the worst-case
(least estimable) direction.
Alternative summary measures include:
\begin{enumerate}[nosep]
    \item \textbf{$D$-efficiency:}
    $r_D = \bigl(\det \FI_{\mathrm{masked}} / \det \FI_{\mathrm{complete}}\bigr)^{1/p}$,
    measuring overall volume.
    \item \textbf{$A$-efficiency:}
    $r_A = \tr(\FI_{\mathrm{complete}}^{-1}) / \tr(\FI_{\mathrm{masked}}^{-1})$,
    measuring average precision.
    \item \textbf{Direction-specific:}
    $r_g = (\nabla g)^\top \FI_{\mathrm{masked}} \nabla g \,/\,
    (\nabla g)^\top \FI_{\mathrm{complete}} \nabla g$
    for a specific function of interest $g(\btheta)$.
\end{enumerate}
Each gives a different perspective.
The minimum-eigenvalue ratio is the most conservative and is
appropriate when all parameters are of equal interest.
\end{remark}

\begin{proposition}[Effective sample size for exponential uniform masking]
\label{prop:ess-expo}
For exponential components under uniform masking with equal rates
$\lambda_1 = \cdots = \lambda_m = \lambda$:
\begin{enumerate}[label=(\alph*)]
    \item $r(\v{\lambda}, w) = \dfrac{m - w}{(m-1)\,w}$ for the minimum-eigenvalue ratio.
    \item $n_{\mathrm{eff}} = n \cdot \dfrac{m - w}{(m-1)\,w}$.
    \item $r_D(\v{\lambda}, w) = \left[\dfrac{m-w}{(m-1)\,w}\right]^{(m-1)/m}$
    for the $D$-efficiency ratio.
\end{enumerate}
In particular, for $m = 3$ and $w = 2$:
$n_{\mathrm{eff}} = n/4$---each masked observation is worth one quarter of a
complete-data observation in the worst-case parameter direction.
\end{proposition}

\begin{proof}
Under equal rates, $\FI_{\mathrm{complete}} = (m\lambda^2)^{-1} I_m$
(from Proposition~\ref{prop:complete-fim} with $\Lambda = m\lambda$).

For $\FI_{\mathrm{masked}}$, by symmetry the matrix has the form
$aI_m + b\v{1}\v{1}^\top$ for constants $a, b$ depending on $m, w, \lambda$.
The eigenvalues are $a + mb$ (eigenvector $\v{1}$) and $a$ (multiplicity
$m - 1$).

To find $a$ and $b$, note that the diagonal entry is
\begin{equation}
    [\FI_{\mathrm{masked}}]_{jj}
    = \frac{1}{\binom{m-1}{w-1} m\lambda}
    \sum_{c \ni j, |c|=w} \frac{1}{w\lambda}
    = \frac{\binom{m-1}{w-1}}{\binom{m-1}{w-1} m w \lambda^2}
    = \frac{1}{mw\lambda^2},
\end{equation}
so $a + b = 1/(mw\lambda^2)$.
The off-diagonal entry (for $j \neq k$) is
\begin{equation}
    [\FI_{\mathrm{masked}}]_{jk}
    = \frac{1}{\binom{m-1}{w-1} m\lambda}
    \sum_{\substack{c \ni j, c \ni k \\ |c|=w}} \frac{1}{w\lambda}
    = \frac{\binom{m-2}{w-2}}{\binom{m-1}{w-1} mw\lambda^2}
    = \frac{w-1}{(m-1)mw\lambda^2},
\end{equation}
so $b = (w-1)/[(m-1)mw\lambda^2]$ and
$a = 1/(mw\lambda^2) - b = (m-w)/[(m-1)mw\lambda^2]$.

The minimum eigenvalue is $a = (m-w)/[(m-1)mw\lambda^2]$.
The minimum eigenvalue of $\FI_{\mathrm{complete}}$ is $1/(m\lambda^2)$.
Hence:
\begin{equation}
    r = \frac{(m-w)/[(m-1)mw\lambda^2]}{1/(m\lambda^2)}
    = \frac{m-w}{(m-1)w}.
\end{equation}
For the special case $w = m - 1$: $r = 1/(m-1)^2$.
For $m = 3$, $w = 2$: $r = 1/4$, so each masked observation
provides $1/4$ the worst-case precision of a complete observation.

For the total rate $\Lambda$ (the $\v{1}$ direction), the retained
information is $a + mb = 1/(mw\lambda^2) + (m-1) \cdot (w-1)/[(m-1)mw\lambda^2]
= [1 + (w-1)]/(mw\lambda^2) = 1/(m\lambda^2)$, which equals the
complete-data information---so the total rate loses nothing to masking,
as expected.
\end{proof}

\begin{remark}[Practical interpretation]
\label{rem:ess-practical}
The effective sample size provides a simple planning tool.
If a practitioner requires the same precision as $n_0$ complete-data
observations, the required number of masked observations is approximately
$n = n_0 / r(\btheta, \pi)$.
For the equal-rate exponential case with $m = 3$, $w = 2$,
$r = 1/4$, so $n = 4 n_0$: four times as many masked observations
are needed to match the precision of complete-data observations in
the worst-case parameter direction.
For the average precision ($A$-efficiency), the factor is milder;
for the total system rate, it is~1 (no penalty).
\end{remark}



%=============================================================================
%  PART III: DESIGN AND DIAGNOSTICS
%=============================================================================

%=============================================================================
\section{Near-Non-Identifiability Diagnostics}
\label{sec:diagnostics}
%=============================================================================

Even when $\rank(\Ctilde) = m$ (formal identifiability), the MLE
can behave poorly if the information about some parameter direction
is very small.
We call this \emph{near-non-identifiability}---the parameter is
technically identifiable but practically unestimable at feasible
sample sizes.

\subsection{Profile likelihood diagnostic}

\begin{definition}[Profile likelihood]
\label{def:profile-likelihood}
For a scalar parameter of interest $\psi = g(\btheta)$, the
\emph{profile likelihood} is
\begin{equation}
    L_p(\psi) = \max_{\btheta: g(\btheta) = \psi} L(\btheta).
\end{equation}
The \emph{profile log-likelihood} is
$\ell_p(\psi) = \max_{\btheta: g(\btheta) = \psi} \ell(\btheta)$.
\end{definition}

The profile likelihood captures all the information about $\psi$
after profiling out the nuisance parameters.
For a well-identified parameter, the profile log-likelihood is
approximately quadratic near its maximum.
Near-non-identifiability manifests as a profile that is
\emph{nearly flat}---the data provide little evidence to distinguish
different values of $\psi$.

\begin{definition}[Profile likelihood curvature ratio]
\label{def:ridge-ratio}
Let $\hat{\psi}$ be the MLE of $\psi$ and let
$\sigma^2_p = -1/\ell_p''(\hat{\psi})$ be the profile-likelihood-based
variance estimate.
The \emph{curvature ratio} for component $j$'s parameters is
\begin{equation}
\label{eq:curvature-ratio}
    \kappa_j = \frac{\lambda_{\min}(\FI_{\mathrm{masked}})}
                    {[\FI_{\mathrm{masked}}^{-1}]_{jj}^{-1}},
\end{equation}
where $[\FI_{\mathrm{masked}}^{-1}]_{jj}^{-1}$ is the reciprocal
of the $j$th diagonal of the inverse Fisher information
(i.e., the marginal information about $\theta_j$ after accounting
for correlations).
\end{definition}

\begin{theorem}[Near-non-identifiability criterion]
\label{thm:near-nonid}
Under standard regularity conditions, component $j$ is
\emph{near-non-identifiable} at sample size $n$ if the
expected profile likelihood curvature satisfies
\begin{equation}
\label{eq:near-nonid}
    n \cdot [\FI_{\mathrm{masked}}^{-1}]_{jj}^{-1}
    < \chi^2_{1,1-\alpha},
\end{equation}
where $\chi^2_{1,1-\alpha}$ is the $(1-\alpha)$-quantile of the
$\chi^2_1$ distribution.
In words: the expected width of the profile-likelihood-based
confidence interval for $\theta_j$ exceeds the parameter space.
\end{theorem}

\begin{proof}[Proof sketch]
The profile log-likelihood ratio statistic for $\psi = \theta_j$ at
the true value is approximately $\chi^2_1$ with expected value 1.
The curvature of the profile log-likelihood at the MLE is approximately
$n / [\FI_{\mathrm{masked}}^{-1}]_{jj}$.
If $n / [\FI_{\mathrm{masked}}^{-1}]_{jj} < \chi^2_{1,1-\alpha}$,
the profile is so flat that the $(1-\alpha)$ confidence interval
extends beyond any meaningful range, indicating that the data
provide essentially no constraint on $\theta_j$.
\end{proof}

\subsection{Eigenvalue diagnostic}

A complementary diagnostic examines the eigenvalues of the observed
(or expected) Fisher information matrix.

\begin{definition}[Condition number diagnostic]
\label{def:condition-diagnostic}
The \emph{identifiability condition number} is
\begin{equation}
    \kappa(\FI_{\mathrm{masked}})
    = \frac{\lambda_{\max}(\FI_{\mathrm{masked}})}
           {\lambda_{\min}(\FI_{\mathrm{masked}})}.
\end{equation}
Large $\kappa$ indicates that some parameter directions are much
harder to estimate than others---the information ``aspect ratio''
is extreme.
\end{definition}

\begin{proposition}[Condition number under equal-rate exponential masking]
\label{prop:condition-expo}
For $m$ exponential components with equal rates $\lambda$ and
uniform masking cardinality $w$:
\begin{equation}
    \kappa(\FI_{\mathrm{masked}})
    = \frac{1 + (w-1) \cdot (m-1)/(m-w) \cdot 1}{1}
    = \frac{w(m-1)}{m - w}.
\end{equation}
This is 1 when $w = 1$ (no masking, isotropic information) and
diverges as $w \to m$ (complete masking, one direction collapses).
For $m = 3$, $w = 2$: $\kappa = 4$.
\end{proposition}

\begin{proof}
From the eigenvalue computations in Proposition~\ref{prop:ess-expo},
$\lambda_{\max} = a + mb$ and $\lambda_{\min} = a$.
Using $a = (m-w)/[(m-1)mw\lambda^2]$ and
$a + mb = 1/(m\lambda^2)$:
\begin{equation}
    \kappa = \frac{1/(m\lambda^2)}{(m-w)/[(m-1)mw\lambda^2]}
    = \frac{(m-1)w}{m-w}.
\end{equation}
\end{proof}

\begin{remark}[Practical thresholds]
\label{rem:thresholds}
We suggest the following guidelines:
\begin{itemize}[nosep]
    \item $\kappa < 10$: mild masking effect, standard MLE reliable.
    \item $10 \leq \kappa < 100$: moderate masking, check profile
    likelihoods and consider regularization.
    \item $\kappa \geq 100$: severe masking, near-non-identifiability
    likely for some parameters; super-component collapse or Bayesian
    methods recommended.
\end{itemize}
These thresholds are heuristic and should be calibrated to the
specific application.
The eigenvalue spectrum of $\FI_{\mathrm{masked}}$ provides a richer
picture than the scalar condition number---a full eigendecomposition
reveals \emph{which} parameter directions are affected.
\end{remark}


%=============================================================================
\section{Optimal Diagnostic Design}
\label{sec:design}
%=============================================================================

The preceding results show that the candidate-set structure determines
how much information is retained.
This raises a natural design question: given a budget constraint on
diagnostic effort, which candidate-set distribution maximizes the
information about $\btheta$?

\subsection{The design space}

Let $\pi = (\pi_c)_{c \in \mathcal{S}}$ denote the probability
distribution over candidate sets, where $\pi_c = \Pr\{\mathcal{C} = c\}$.
Under C1--C2--C3, the analyst can influence $\pi$ by choosing the
diagnostic procedure.
More precise diagnostics produce smaller candidate sets
(more singletons), but are typically more expensive.

We model the diagnostic cost as a function of the candidate-set
distribution: $\text{Cost}(\pi) = \sum_c \pi_c \cdot \text{cost}(c)$,
where $\text{cost}(c)$ is the per-observation diagnostic cost for
resolving the failure cause to the candidate set~$c$.
A natural assumption is that $\text{cost}(c)$ is decreasing in $|c|$
(smaller sets require more diagnostic effort).

\subsection{D-optimal design}

\begin{definition}[$D$-optimal diagnostic design]
\label{def:d-optimal}
The \emph{$D$-optimal diagnostic design} maximizes the determinant
of the Fisher information matrix subject to a cost constraint:
\begin{equation}
\label{eq:d-optimal}
    \pi^* = \arg\max_\pi \det \FI_{\mathrm{masked}}(\btheta; \pi)
    \quad \text{subject to} \quad
    \text{Cost}(\pi) \leq B,
\end{equation}
where $B$ is the diagnostic budget and
$\FI_{\mathrm{masked}}(\btheta; \pi)$ denotes the Fisher information
under masking distribution $\pi$.
\end{definition}

The $D$-optimality criterion minimizes the volume of the asymptotic
confidence ellipsoid for $\btheta$
\citep{atkinson2007,pukelsheim2006}.

\begin{proposition}[$D$-optimal design structure]
\label{prop:d-optimal-structure}
For the diagnostic design problem \eqref{eq:d-optimal}:
\begin{enumerate}[label=(\alph*)]
    \item If singleton candidate sets $\{j\}$ are available for all
    $j$ (i.e., exact identification is always possible), the
    $D$-optimal design with no budget constraint assigns all
    probability to singletons, recovering $\FI_{\mathrm{complete}}$.
    \item Under the cost model
    $\text{cost}(c) = \alpha / |c|$ (cost proportional to diagnostic
    precision), the $D$-optimal design is a mixture of singleton and
    non-singleton candidate sets, with the mixture proportions
    determined by the Lagrange conditions.
    \item For exponential components with equal rates, the $D$-optimal
    design among uniform masking configurations is $w = 1$ (no masking)
    whenever the budget allows, and $w = m - 1$ (maximum masking)
    when the budget forces the largest possible candidate sets.
\end{enumerate}
\end{proposition}

\begin{proof}[Proof sketch]
Part~(a) follows because $\FI_{\mathrm{complete}} \succeq
\FI_{\mathrm{masked}}$ (Theorem~\ref{thm:fi-decomposition}), and
$\det(\cdot)$ is monotone with respect to the Loewner order on
positive definite matrices.

Part~(b) is a standard convex optimization result: the Fisher
information is a linear function of $\pi$ (each candidate set
contributes additively to the information), and the cost constraint
is linear.
The optimal $\pi$ is found at a vertex or on the boundary of the
feasible simplex, which corresponds to a mixture of at most $m + 1$
support points (by Carath\'eodory's theorem applied to the information
matrices).

Part~(c) follows from the symmetry of the equal-rate problem:
among uniform masking configurations with fixed $w$, the information
matrix depends only on $w$, and
$\det \FI_{\mathrm{masked}}(\v{\lambda} \mid w)$ is decreasing
in $w$ (by Theorem~\ref{thm:fi-monotone} and the fact that the
determinant is monotone under Loewner ordering).
\end{proof}

\begin{remark}[Tradeoff between resolution and sample size]
\label{rem:design-tradeoff}
In practice, the budget constraint often means choosing between
fewer high-precision observations (small candidate sets) and more
low-precision observations (large candidate sets).
The effective sample size framework
(Section~\ref{sec:ess}) makes this tradeoff explicit:
$n_1$ observations with retention ratio $r_1$ provide the same
information as $n_2$ observations with retention ratio $r_2$ when
$n_1 r_1 = n_2 r_2$.
If a singleton observation costs $\alpha$ and a size-$w$ observation
costs $\alpha/w$, the budget $B$ buys either $B/\alpha$ singletons
or $Bw/\alpha$ masked observations.
The design decision reduces to whether
$(B/\alpha) \cdot 1 \gtrless (Bw/\alpha) \cdot r(\btheta, w)$,
i.e., whether $r(\btheta, w) \cdot w \gtrless 1$.
For the equal-rate exponential case,
$r \cdot w = (m - w)/[(m-1)] < 1$ for $w \geq 2$, so singletons
are always preferred under this cost model.
Under different cost structures (e.g., fixed cost per observation
plus variable diagnostic cost), the optimal design may favor a
mixture of masking levels.
\end{remark}


%=============================================================================
\section{Connection to Missing Data Theory}
\label{sec:missing-data}
%=============================================================================

The masking of failure causes fits naturally within the missing data
framework of \citet{little2002}.
This connection provides additional insight and links our Fisher
information decomposition to established missing-data diagnostics.

\subsection{Masking as a missing data mechanism}

In the complete-data formulation, an observation is $(T_i, K_i)$.
Masking replaces the exact cause $K_i$ with the coarser candidate
set $\mathcal{C}_i$.
This is a form of \emph{coarsening}: the observed data is a
many-to-one function of the complete data.

Under C1--C2--C3, the masking mechanism is \emph{ignorable} in the
sense of \citet{little2002}:
\begin{itemize}[nosep]
    \item C1 ensures that the candidate set is consistent with the
    complete data (the missing-data pattern respects the truth).
    \item C2 and C3 together ensure that the probability of the
    observed masking pattern, given the complete data, does not depend
    on the parameter $\btheta$---this is the missing-at-random (MAR)
    condition, combined with distinctness of parameters.
\end{itemize}
Under ignorability, the observed-data likelihood is proportional to
the complete-data likelihood marginalized over the unobserved
component, which is exactly the C1--C2--C3 likelihood
\eqref{eq:likelihood-contribution}.

\subsection{Fraction of missing information}

\begin{definition}[Fraction of missing information]
\label{def:fmi}
The \emph{fraction of missing information} about the $j$th parameter
is
\begin{equation}
\label{eq:fmi}
    \gamma_j
    = 1 - \frac{[\FI_{\mathrm{masked}}^{-1}]_{jj}^{-1}}
               {[\FI_{\mathrm{complete}}^{-1}]_{jj}^{-1}}
    = 1 - \frac{\text{marginal masked info about } \theta_j}
               {\text{marginal complete info about } \theta_j}.
\end{equation}
The matrix-valued fraction of missing information is
\begin{equation}
    \Gamma = I_p - \FI_{\mathrm{complete}}^{-1} \FI_{\mathrm{masked}},
\end{equation}
whose eigenvalues lie in $[0, 1]$.
\end{definition}

\begin{proposition}[FMI for exponential uniform masking]
\label{prop:fmi-expo}
For exponential components under uniform masking with equal rates:
\begin{enumerate}[label=(\alph*)]
    \item The eigenvalues of $\Gamma$ are 0 (multiplicity 1,
    corresponding to the total rate $\Lambda$) and
    $1 - (m-w)/[(m-1)w] \cdot (m-1) = w/(m-1) \cdot (m-1-m+w)/(w) = \ldots$

    More precisely, $\Gamma = I - \FI_{\mathrm{complete}}^{-1}
    \FI_{\mathrm{masked}}$.
    With $\FI_{\mathrm{complete}} = (m\lambda^2)^{-1} I$ and
    $\FI_{\mathrm{masked}} = aI + b\v{1}\v{1}^\top$:
    $\Gamma = I - m\lambda^2 (aI + b\v{1}\v{1}^\top)
    = (1 - m\lambda^2 a)I - m\lambda^2 b \v{1}\v{1}^\top$.
    The eigenvalues of $\Gamma$ are $1 - m\lambda^2(a + mb) = 0$
    (eigenvector $\v{1}$) and $1 - m\lambda^2 a$
    (multiplicity $m - 1$).
    Using $a = (m-w)/[(m-1)mw\lambda^2]$:
    $1 - m\lambda^2 \cdot (m-w)/[(m-1)mw\lambda^2]
    = 1 - (m-w)/[(m-1)w]$.

    \item For $m = 3$, $w = 2$:
    $\gamma = 1 - 1/2 = 1/2$ for each individual rate, and
    $\gamma = 0$ for the total rate.
    Half the information about each individual rate is lost to masking.
\end{enumerate}
\end{proposition}

\begin{remark}[Connection to EM convergence rate]
\label{rem:em-convergence}
The fraction of missing information $\gamma$ also governs the
convergence rate of the EM algorithm \citep{dempster1977}.
For the masked series system, the EM algorithm treats $K_i$ as the
latent variable and iterates between computing
$\E[K_i \mid T_i, \mathcal{C}_i; \btheta^{(t)}]$ (E-step) and
maximizing the expected complete-data log-likelihood (M-step).
The convergence rate is approximately
$(1 - \gamma_{\min})^{-1}$ iterations per digit of accuracy, where
$\gamma_{\min}$ is the largest eigenvalue of $\Gamma$
\citep{dempster1977}.
High fractions of missing information (heavy masking) slow EM
convergence, another manifestation of the same information loss.
\end{remark}



%=============================================================================
%  PART IV: VALIDATION AND DISCUSSION
%=============================================================================

%=============================================================================
\section{Simulation Study}
\label{sec:simulation}
%=============================================================================

We design a comprehensive simulation study to validate the theoretical
results and explore their finite-sample behavior.
The study systematically varies the factors that govern identifiability
and information loss.

\subsection{Experimental design}

\paragraph{Factors.}
The simulation varies five factors in a factorial design:

\begin{table}[ht]
\centering
\caption{Simulation study factors and levels.}
\label{tab:simulation-factors}
\begin{tabular}{@{}lll@{}}
\toprule
\textbf{Factor} & \textbf{Levels} & \textbf{Purpose} \\
\midrule
Number of components $m$ & 3, 5, 10 & Scalability \\
Masking cardinality $w$ & 1, $\lceil m/2 \rceil$, $m-1$
    & Information loss gradient \\
Sample size $n$ & 50, 200, 1000, 5000 & Asymptotic convergence \\
Rate configuration & Equal, moderate (1:2:3),
    & Parameter asymmetry \\
    & strong (1:3:10) & \\
Distribution family & Exponential, Weibull($k{=}2$)
    & Shape sensitivity \\
\bottomrule
\end{tabular}
\end{table}

\paragraph{Candidate-set structures.}
For each $(m, w)$ combination, we consider three candidate-set
structures:
\begin{enumerate}[nosep]
    \item \textbf{Uniform random:} Each candidate set of size $w$
    containing the true cause is equally likely (the uniform masking
    model).
    \item \textbf{Block structure:} Components are grouped into
    blocks, and the candidate set is the block containing the true
    cause (fixed structure, varying block sizes).
    \item \textbf{Hierarchical:} Components form a tree; the
    candidate set is the subtree at a random level containing the
    true cause.
\end{enumerate}

\paragraph{Response variables.}
For each configuration, we generate $R = 10{,}000$ replicate data
sets and compute:
\begin{enumerate}[nosep]
    \item \textbf{MLE bias and MSE} for each component parameter.
    \item \textbf{Empirical coverage} of nominal 95\% Wald confidence
    intervals based on the Fisher information.
    \item \textbf{Eigenvalue spectrum} of the observed Fisher
    information, compared to the theoretical $\FI_{\mathrm{masked}}$.
    \item \textbf{Effective sample size} $n_{\mathrm{eff}}$ (empirical
    vs.\ theoretical).
    \item \textbf{Fraction of missing information} $\gamma_j$
    (empirical vs.\ theoretical).
    \item \textbf{Convergence rate} of the numerical optimizer
    (iterations to convergence).
    \item \textbf{Non-identifiability rate}: fraction of replicates
    where the optimizer fails to converge or produces boundary estimates.
\end{enumerate}

\subsection{Theoretical predictions}

The simulation is designed to test the following predictions from the
theory developed in Sections~\ref{sec:fi-decomposition}--\ref{sec:ess}:

\begin{enumerate}[label=(P\arabic*)]
    \item \textbf{MSE scaling:} Empirical MSE should scale as
    $1/(n \cdot \lambda_{\min}(\FI_{\mathrm{masked}}))$, with the
    ratio MSE$_{\mathrm{masked}}$ / MSE$_{\mathrm{complete}}$
    converging to $1/r(\btheta, w)$.
    \item \textbf{Coverage:} Wald interval coverage should approach
    95\% as $n$ increases, with slower convergence for larger $w$
    and more asymmetric rate configurations.
    \item \textbf{Monotonicity:} Information loss metrics should
    increase monotonically with $w$, as predicted by
    Theorem~\ref{thm:fi-monotone}.
    \item \textbf{Effective sample size:} Empirical
    $n_{\mathrm{eff}}$ should match the theoretical prediction
    from Definition~\ref{def:ess}.
    \item \textbf{Condition number:} The eigenvalue ratio of the
    empirical Fisher information should converge to
    $\kappa(\FI_{\mathrm{masked}})$ from
    Proposition~\ref{prop:condition-expo}.
    \item \textbf{Non-identifiability rates:} The fraction of
    non-identifiable replicates should decrease as $n$ increases
    and increase with $w$ and rate asymmetry, consistent with
    the near-non-identifiability analysis of Section~\ref{sec:diagnostics}.
\end{enumerate}

\subsection{Implementation}

The simulation will be implemented using the
\texttt{maskedcauses}~R package \citep{towell2025maskedcauses}
for data generation and the
\texttt{algebraic.mle}~R package \citep{towell2023algebraicmle}
for MLE computation and Fisher information analysis.
For exponential components with $w = m - 1$, the closed-form MLE
from \citet{towell2026exponential} is used directly.
For all other configurations, L-BFGS-B optimization
\citep{byrd1995} is used with multiple random starting points.

\paragraph{Note.}
\emph{The simulation results are to be computed.
Tables and figures will be populated once the simulation code is
executed.
The experimental design above specifies the complete protocol;
the theoretical predictions (P1)--(P6) serve as the hypotheses
to be tested.}


%=============================================================================
\section{Discussion}
\label{sec:discussion}
%=============================================================================

\subsection{Summary of contributions}

This paper develops the quantitative theory of identifiability and
information loss for masked series systems under the C1--C2--C3
framework.
The main contributions are:

\begin{enumerate}
    \item \textbf{Confounding graph}
    (Section~\ref{sec:graph}): A graph-theoretic characterization of
    partial identifiability.
    The connected components of the confounding graph are the
    super-components---the minimal units whose aggregate hazards are
    identifiable.
    This gives practitioners a visual and computational tool for
    understanding which parameters can be separated from which.

    \item \textbf{Matrix rank condition}
    (Section~\ref{sec:rank-condition}): For exponential components,
    the identifiability condition is exactly
    $\rank(\Ctilde) = m$---both necessary and sufficient.
    This is a self-contained linear algebra result that can be checked
    computationally from the candidate-set structure alone.

    \item \textbf{Fisher information decomposition}
    (Sections~\ref{sec:fi-decomposition}--\ref{sec:expo-info}):
    The masked-data Fisher information decomposes as
    $\FI_{\mathrm{masked}} = \FI_{\mathrm{complete}} - \Delta$,
    with $\Delta \succeq 0$.
    For exponential components under uniform masking, all terms
    have closed forms.

    \item \textbf{Monotone information loss}
    (Section~\ref{sec:monotone}): The information loss increases
    monotonically with the masking cardinality, with uniform masking
    at $w = m - 1$ providing the worst-case benchmark.

    \item \textbf{Effective sample size}
    (Section~\ref{sec:ess}): A scalar summary that translates the
    matrix-valued information loss into practical terms: how many
    masked observations equal one complete observation?

    \item \textbf{Diagnostics and design}
    (Sections~\ref{sec:diagnostics}--\ref{sec:design}): Profile
    likelihood diagnostics for detecting near-non-identifiability,
    and $D$-optimal diagnostic design for maximizing information
    under cost constraints.

    \item \textbf{Missing data connection}
    (Section~\ref{sec:missing-data}): The fraction of missing
    information links the Fisher decomposition to established
    missing data theory and EM convergence rates.
\end{enumerate}

\subsection{Mutual information versus Fisher information}

Two notions of ``information about the failure cause'' arise in
this paper:
\begin{itemize}[nosep]
    \item \textbf{Mutual information} $\MI(K; \mathcal{C})$:
    how much the candidate set reveals about \emph{which component
    failed} (a combinatorial quantity, independent of $\btheta$).
    \item \textbf{Fisher information} $\FI_{\mathrm{masked}}(\btheta)$:
    how precisely the \emph{parameter vector} $\btheta$ can be
    estimated from the masked data (a statistical quantity, dependent
    on $\btheta$).
\end{itemize}

These are related but distinct.
The mutual information depends only on the masking mechanism and the
marginal distribution of $K$, not on the parametric family.
The Fisher information depends on the full likelihood, including the
functional form of the hazards.
Monotonicity in $w$ holds for both, but the mutual information is
a scalar while the Fisher information is a matrix, and their orderings
can differ in specific parameter directions.

A deeper connection between the two---for instance, conditions under
which higher $\MI(K; \mathcal{C})$ implies higher
$\FI_{\mathrm{masked}}$ in a matrix sense---remains an open problem.
For the equal-rate exponential case, the relationship is clean:
$\MI(K; \mathcal{C}_w) = \ln(m/w)$ while the information retention
ratio scales as $(m-w)/[(m-1)w]$, and both are monotonically
decreasing in $w$.
For unequal rates and non-exponential families, the relationship
is more complex.

\subsection{Limitations and open problems}

\begin{enumerate}
    \item \textbf{Beyond uniform masking.}
    The closed-form results in Sections~\ref{sec:expo-info}
    and~\ref{sec:monotone} assume uniform masking.
    Under general C1--C2--C3 masking, the Fisher information depends
    on the unknown masking distribution $\pi$, which is profiled out
    of the likelihood (by C3) but affects the information matrix.
    Extending the decomposition to general C2 models is an important
    direction.

    \item \textbf{Non-exponential families.}
    The Fisher information decomposition
    (Theorem~\ref{thm:fi-decomposition}) holds for any parametric
    family, but the closed-form expressions are specific to exponential
    components.
    For Weibull and other families, numerical computation of the Fisher
    information is required.
    Developing closed-form or asymptotic expressions for the
    masking-loss matrix under Weibull components is the subject of
    ongoing work \citep{towell2026weibull}.

    \item \textbf{General parametric family extension.}
    For families with time-varying hazards (e.g., Weibull with
    distinct shape parameters), the identifiability condition can be
    weaker than $\rank(\Ctilde) = m$---the functional diversity of
    the hazards provides additional constraints
    (Remark~\ref{rem:expo-vs-general}).
    Characterizing the exact identifiability condition for general
    families is open.

    \item \textbf{Optimal design beyond equal rates.}
    The $D$-optimal design results
    (Section~\ref{sec:design}) are developed for the symmetric
    (equal-rate) case.
    Under unequal rates, the optimal design depends on the unknown
    $\btheta$, leading to \emph{locally optimal} designs
    \citep{chernoff1953}.
    Sequential or Bayesian design procedures that adapt the diagnostic
    strategy as data accumulate are a natural extension.

    \item \textbf{Minimax candidate-set distribution.}
    What candidate-set distribution minimizes the worst-case
    information loss over all $\btheta$ in some class?
    This minimax formulation connects to robust experimental design
    and may yield qualitatively different solutions from
    $D$-optimality.
\end{enumerate}

\subsection{Connections to the wider ecosystem}

This paper is the primary companion to the foundation paper
\citep{towell2025masked}, extending its binary identifiability theorem
into a quantitative framework.
The exponential Fisher information results re-derived here in
self-contained form were first established in
\citet{towell2026exponential}; the present paper places them in
the broader context of general parametric families and information
decomposition.

The effective sample size and information loss results have direct
implications for the model selection paper
\citep{towell2026model-selection}: the power of likelihood ratio
tests for nested models (e.g., homogeneous vs.\ heterogeneous Weibull)
depends on the Fisher information under masking, and the
near-non-identifiability diagnostics developed here can guide the
choice of model complexity.

The observation composition paper (companion~03) establishes
closure properties of C1--C2--C3 under data combination; the
information decomposition here provides the quantitative complement,
showing how the combined Fisher information relates to the component
information matrices.



%=============================================================================
\section{Conclusion}
\label{sec:conclusion}
%=============================================================================

The foundation paper's identifiability theorem tells us \emph{whether}
component parameters can be recovered from masked series system data.
This paper answers the harder questions: \emph{how precisely}, and
\emph{what to do about it}.

The confounding graph provides a structural diagnostic: its connected
components are the irreducible units of identifiability, and their
count tells the practitioner exactly how many parameters can be
separated.
The Fisher information decomposition provides a quantitative
diagnostic: the masking-loss matrix $\Delta$ measures the price of
incomplete diagnostics in units of statistical precision.
The effective sample size translates this into the currency
practitioners understand: how many more observations do I need?

For exponential components under uniform masking, the entire theory
admits closed forms.
The information retention ratio $(m-w)/[(m-1)w]$ gives a clean rule
of thumb: with $m = 5$ components and $w = 4$ masking cardinality,
each masked observation is worth $1/12$ of a complete observation in
the worst-case parameter direction.
Twenty-five percent diagnostic improvement (reducing $w$ from 4 to 3)
roughly doubles the effective sample size.

The practical tools---profile likelihood diagnostics, the condition
number test, $D$-optimal design---move the theory from mathematical
characterization to engineering decision-making.
When the condition number of the Fisher information exceeds a
threshold, the analyst knows to collapse super-components or
regularize.
When designing a new diagnostic procedure, the $D$-optimality
framework provides the principled tradeoff between diagnostic
precision and sample size.

The connection to missing data theory situates the entire development
within the established statistical framework of \citet{little2002},
providing both theoretical grounding (the fraction of missing
information as the eigenvalues of $\Gamma$) and computational
guidance (the EM convergence rate as a function of $\gamma$).

Open problems remain---notably the extension to non-exponential
families in closed form, the minimax diagnostic design, and the
deeper connection between mutual information and Fisher
information---but the framework developed here provides the
foundation for addressing them.


%=============================================================================
% Bibliography
%=============================================================================

\bibliographystyle{plainnat}
\bibliography{refs}

\end{document}
